\subsection{Ebenendarstellungen umformen}
\label{subsec:analytic_geometry_converting_plane_representations}

Parameter- und Normalenform beschreiben dieselbe Ebene aus zwei Blickwinkeln.
Die Umrechnung macht jeweils die für eine Aufgabe günstigere Information
sichtbar.

\begin{center}
  \CachedTikzFigure{content/resources/chapters/analytic_geometry/tikz/plane_representation_conversion.tex}
\end{center}

\begin{KeyPointBox}{Wechsel zwischen Parameter- und Normalenform}{analytic_geometry_plane_representation_conversion}
  \[
    \begin{array}{c@{\qquad\Longrightarrow\qquad}c}
      \text{Parameterform} & \text{Normalenform}\\[0.3em]
      \vec x=\vec p+\lambda\vec a+\mu\vec b
      & (\vec x-\vec p)\cdot(\vec a\times\vec b)=0
    \end{array}
  \]
  In der Gegenrichtung wählt man einen Punkt $P$ der Ebene und zwei linear
  unabhängige Lösungen von $\vec n\cdot\vec r=0$ als Richtungsvektoren.
\end{KeyPointBox}

\begin{ExampleBox}{Parameterform in Normalenform überführen}{analytic_geometry_parameter_to_normal_example}
  Gegeben ist die bereits bekannte Ebene
  \[
    E:\quad
    \vec x=\begin{pmatrix}6\\3\\-5\end{pmatrix}
      +\lambda\begin{pmatrix}1\\1\\-2\end{pmatrix}
      +\mu\begin{pmatrix}-1\\3\\1\end{pmatrix}.
  \]
  Ein Normalenvektor entsteht aus dem Kreuzprodukt der Richtungsvektoren:
  \begin{ClozeArea}[34]{12}
    \[
      \vec n
      =\begin{pmatrix}1\\1\\-2\end{pmatrix}
        \times\begin{pmatrix}-1\\3\\1\end{pmatrix}
      =\begin{pmatrix}
          1\cdot1-(-2)\cdot3\\
          (-2)\cdot(-1)-1\cdot1\\
          1\cdot3-1\cdot(-1)
        \end{pmatrix}
      =\begin{pmatrix}7\\1\\4\end{pmatrix}.
    \]
  \end{ClozeArea}
  Damit lauten Normalen- und Koordinatenform
  \begin{ClozeArea}[34]{8}
    \[
      \begin{aligned}
        E:\quad
        \left(\vec x-\begin{pmatrix}6\\3\\-5\end{pmatrix}\right)
          \cdot\begin{pmatrix}7\\1\\4\end{pmatrix}&=0,\\[0.8em]
        7x_1+x_2+4x_3&=25.
      \end{aligned}
    \]
  \end{ClozeArea}
\end{ExampleBox}

\begin{ExampleBox}{Normalenform in Parameterform überführen}{analytic_geometry_normal_to_parameter_example}
  Gegeben ist
  \[
    E:\quad
    \left(\vec x-\begin{pmatrix}2\\1\\0\end{pmatrix}\right)
      \cdot\begin{pmatrix}2\\-1\\1\end{pmatrix}=0.
  \]
  Der Aufpunkt bleibt erhalten. Für Richtungsvektoren $\vec r$ muss
  $\vec n\cdot\vec r=0$ gelten. Zwei unabhängige Lösungen sind
  \begin{ClozeArea}[34]{8}
    \[
      \vec a=\begin{pmatrix}1\\2\\0\end{pmatrix},
      \qquad
      \vec b=\begin{pmatrix}1\\1\\-1\end{pmatrix},
      \qquad
      \vec n\cdot\vec a=\vec n\cdot\vec b=0.
    \]
  \end{ClozeArea}
  Somit besitzt die Ebene beispielsweise die Parameterdarstellung
  \[
    E:\quad
    \vec x=\begin{pmatrix}2\\1\\0\end{pmatrix}
      +\lambda\begin{pmatrix}1\\2\\0\end{pmatrix}
      +\mu\begin{pmatrix}1\\1\\-1\end{pmatrix},
    \qquad \lambda,\mu\in\R.
  \]
\end{ExampleBox}
